% Source: https://artofproblemsolving.com/community/c3h1342537_an_incorrect_proof_that_sqrt2__1
% Stable identity: aops-topic-1342537-q01
There are two errors.

\textbf{1. The main error.} Taking a square root means raising to the power $\dfrac{1}{2}$: $\sqrt{x^2}=\left(x^2\right)^{\dfrac{1}{2}}$. By the power rule $\left(x^a\right)^b=x^{ab}$, this is $x^{2\cdot\dfrac{1}{2}}=x^1$. The proof instead writes $x^{\left(2^{\dfrac{1}{2}}\right)}=x^{\sqrt2}$, raising the \emph{exponent} $2$ to the power $\dfrac{1}{2}$. But $\left(x^2\right)^{\dfrac{1}{2}}$ and $x^{\left(2^{\dfrac{1}{2}}\right)}$ are different numbers; for example with $x=4$: $\left(4^2\right)^{\dfrac{1}{2}}=4$ while $4^{\sqrt2}\approx7.1$. Exponents of a power multiply, they are not stacked: $(a^b)^c=a^{bc}\ne a^{\left(b^c\right)}$ in general.

\textbf{2. A smaller error.} Even the first line is not right for all real $x$: $\sqrt{x^2}=|x|$, which equals $x$ only when $x\ge0$.

Finally, even if $x^1=x^{\sqrt2}$ held for some $x$, one could not conclude $1=\sqrt2$ when $x=0$ or $x=1$.
